% LSPV-3 REPORT::CYCLEGAIN1::2026-10-04T19:37:00-07:00
% !TEX program = xelatex
\documentclass[11pt]{article}
\usepackage[margin=.86in]{geometry}
\usepackage{amsmath,amssymb}
\begin{document}
\section*{LSPV-3 CYCLEGAIN1: cycle-level contraction with same-scale typed edges}
\textbf{GLOBAL STATUS: OPEN.} P versus NP remains open. This report refines CYCLECLOSE1. It proves broader conditional recurrence accounting in which some typed recursive edges may preserve input scale, provided same-scale recursion is source-derived and well founded and every directed return cycle contains genuine contraction. It does not construct a SAT shrinker, close TASK-001, force arbitrary SAT algorithms into the typed graph, or supply a nonrelativizing/nonalgebrizing hinge.

\section*{CORPUS REVIEWED}
AGENTS.md, CODEX.md, the pnp-research skill, root main.tex, verified MAIN1854, RECURCLOSE1, SELECTCHARGE1, MULTICHANNEL1, CYCLECLOSE1, AFFQUAD1, current TASK-001 search results, the causal TASK-004/TASK-005 state represented in root and prior synthesis, TASK-008 search results, and independent hostile-control reports were surveyed. Repository search found no verified CYCLEGAIN1 artifact. AFFQUAD1 is the newest verified LSPV-2 artifact materially relevant here and leaves the semantic return bridge OPEN.

\section*{Predecessor recovery check}
The immediately preceding specialist is LSPV-2. AFFQUAD1 exists, is substantive, and has the required header and final sentinel. No predecessor recovery request is required.

\section*{Primary literature preflight}
Buss--Williams supplies the alternation-trading SAT time-space frontier only after a model-matched time/space upper bound is proved. Fortnow--Lipton--van Melkebeek--Viglas gives unconditional SAT time-space lower bounds in its stated random-access model. General recursion literature studies mutually recursive call graphs, but no external recurrence theorem is imported as a SAT bridge. CYCLEGAIN1 is proved directly.

\section*{Exact typed model}
Fix finitely many deterministic exact procedure types \(P_1,\ldots,P_m\). Type \(i\), on an input of length \(n\), performs charged local work \(n^{r_i+o(1)}\) and may use finitely many channels \(e:i\to j\), with at most \(n^{\eta_e+o(1)}\) calls on explicitly constructed target inputs.

A fractional edge has \(0\le\theta_e<1\) and target length at most \(n^{\theta_e+o(1)}\). A same-scale edge has \(\theta_e=1\), target length at most \(n\), and strictly decreases a source-derived finite rank \(h(i,x)\in\{0,1,\ldots,H\}\), where \(H\) is independent of \(n\). In the finite type-only form below, the same-scale subgraph is acyclic. This condition is not free and must be proved from the actual source construction.

Acquisition, normalization, census/output, amplification, indexing, compilation, route work, selector/query construction, verification, writing, transport, rereading, recomputation, recursion-stack space, and live buffers are separately charged. SAME\_REALIZATION requires recursive substitution to answer exactly the declared source calls.

\section*{PROVED/EXACT: cycle-gain closure}
Let \(C_i\ge r_i\) and for every edge \(e:i\to j\), \(C_i\ge \eta_e+\theta_e C_j\).

\textbf{Theorem CYCLEGAIN1.} Assume \(0\le\theta_e\le1\) for all edges and every directed cycle \(\gamma\) satisfies
\[
\boxed{\Theta_\gamma:=\prod_{e\in\gamma}\theta_e<1.}
\]
Assume scale-one calls are nonexpanding in actual encoded length and well founded as above. Then the max-affine exponent system has a finite least solution \(C\), and depth-first recursive substitution yields
\[
\boxed{T_i(n)\le n^{C_i+o(1)}}
\]
for every type \(i\), independently of the original exponent \(K\).

\textbf{Proof.} Since all \(\theta_e\le1\), a cycle has product one exactly when all its edges have scale one. Thus the hypothesis is equivalent to acyclicity of the scale-one subgraph. Let
\[
\lambda=\max\{\theta_e:\theta_e<1\}<1.
\]
Along any path, at most \(m-1\) consecutive scale-one edges occur before termination or a fractional edge. Hence every block of \(m\) recursive edges contains a factor at most \(\lambda\), so
\[
\prod_{q=1}^{t}\theta_{e_q}\le\lambda^{\lfloor t/m\rfloor}.
\]
With \(\eta_*=\max_e\eta_e\) and \(r_*=\max_i r_i\), weighted multiplicity along any path is bounded by
\[
m\eta_*\sum_{q\ge0}\lambda^q+r_*<\infty.
\]
The supremum of exact path weights from each start type therefore gives a finite feasible exponent vector, and monotone iteration from local work gives the least one. Same-scale recursion terminates by its finite rank; fractional edges eventually reduce actual input length below a fixed base threshold. Strong induction on input length together with the same-scale rank proves the bound. \(\square\)

\section*{PROVED/EXACT: cycle inequality}
For a directed cycle \(\gamma=(e_1,\ldots,e_d)\), define
\[
H_\gamma=\eta_{e_1}+\theta_{e_1}\eta_{e_2}+\cdots+
\left(\prod_{q=1}^{d-1}\theta_{e_q}\right)\eta_{e_d}.
\]
Composition gives \(C_i\ge H_\gamma+\Theta_\gamma C_i\). If \(\Theta_\gamma<1\),
\[
\boxed{C_i\ge\frac{H_\gamma}{1-\Theta_\gamma}}.
\]
If \(\Theta_\gamma=1\) and \(H_\gamma>0\), no finite exponent vector satisfies that cycle. This is a failure of this exponent-erasure construction, not a SAT lower bound.

\section*{Correction to CYCLECLOSE1}
CYCLECLOSE1 remains valid under its stronger sufficient condition \(\theta_e<1\) on every edge. The blanket typed-recursion claim that one \(\theta_e=1\) edge necessarily preserves \(K\) is \textbf{FALSIFIED/RETIRED}. A same-scale edge may be absorbed inside a source-matched well-founded typed system whose complete return cycles contract.

For example,
\[
C_1\ge a+C_2,\qquad C_2\ge b+\tfrac12 C_1
\]
has finite requirements \(C_1\ge2(a+b)\) and \(C_2\ge a+2b\). The first edge is same-scale, while the complete return cycle has gain \(1/2\).

SELECTCHARGE1 remains correct for an unresolved direct invocation of the original black-box SAT decider \(A\) on current-scale inputs. Such a call still contributes the unknown exponent \(K\). CYCLEGAIN1 applies only after that call has itself been replaced by the proved source-matched typed system.

\section*{PROVED/EXACT: space gate}
Time contraction does not imply low space. If a same-scale child query of length \(\Theta(n)\) is explicitly materialized, that representation uses
\[
\boxed{S(n)\ge n^{1-o(1)}}.
\]
Under explicit query buffers, \(\rho\ge\max_e\theta_e\). Thus a materialized \(\theta_e=1\) edge prevents direct use of the subpolynomial-space SAT frontier even when time cycles contract. A same-scale route therefore needs a separately proved streaming, implicit, or compact query interface whose construction, access, evaluation, and retained workspace are fully charged.

\section*{Conservative reconstruction and quantifier ledger}
The standing reconstruction predicate remains
\[
\Psi(x,Q)\equiv\exists\sigma\;\forall u\;\exists z\;\forall y\;R(x,Q,\sigma,u,z,y).
\]
No quantifier is commuted or canonically bound by CYCLEGAIN1. The conservative ledger remains
\[
D(K)=\max\{a,q+bK\},\quad R(K)=\ell+c_0K^2,\quad F(K)=\max\{f,c_FK\},
\]
\[
E(K)=\min\{D(K),R(K),F(K)\},\qquad T_{\rm total}(N)\le N^{\max\{G,E(K)\}+o(1)}.
\]
CYCLEGAIN1 applies only after a source-matched return construction exposes the declared typed channels.

\section*{Hostile controls}
DELAY does not create contraction. COPY increases multiplicity charge. RECOMP belongs in local work or explicit edges. INTERP changes the graph and requires a fresh full ledger. Padding cannot manufacture fractional scale when measured against the unpadded input. Compact representations evade the explicit-buffer space gate only if acquisition and evaluation are proved and charged. Horn, dual-Horn, 2-SAT, affine/Tseitin, bounded treewidth, and independent products remain easy controls and do not imply unrestricted SAT hardness.

\section*{Barrier and transfer audit}
CYCLEGAIN1 relativizes under consistent oracle replacement and is not the missing nonrelativizing hinge. It is combinatorial recurrence accounting, but that alone does not make the final route nonalgebrizing. Buss--Williams and Fortnow--Lipton--van Melkebeek--Viglas transfer only after a matching low-space upper bound. Call-graph and small-gain language are structural analogies, not SAT source/acquisition theorems.

\section*{DUPLICATION / CONTRADICTION AUDIT}
RECURCLOSE1 is the one-type one-channel case. MULTICHANNEL1 is the one-type finite-channel case with fractional channels. CYCLECLOSE1 is finite typed recursion under per-edge fractional contraction. CYCLEGAIN1 relaxes this to cycle-level contraction plus well-founded same-scale recursion. MAIN1854 and AFFQUAD1 say any \(\theta=1\) channel blocks exponent erasure. That remains correct for an unresolved direct black-box SAT channel, but is stale if read as a theorem about an already source-matched finite typed system.

\section*{WHAT CHANGED RELATIVE TO ROOT MAIN}
Root main currently treats a directed cycle containing a current-scale dependency as obstructing. CYCLEGAIN1 sharpens this: a cycle may contain scale-one edges and still close if its total product gain is below one. The typed obstruction is a noncontracting cycle, plus the separate space cost of any materialized same-scale query.

\section*{ADOPTED / REJECTED / RETIRED PRIOR COMPONENTS}
\textbf{Adopted:} RECURCLOSE1, MULTICHANNEL1, CYCLECLOSE1, SELECTCHARGE1 for unresolved direct calls to \(A\), FULLSKOLEM1, the conservative ledger, strict charge separation, and SAME\_REALIZATION.
\textbf{Rejected as sufficient:} feedback alone, algebraic/product-span collisions, compact representations without acquisition, and type relabeling that hides a direct call to \(A\).
\textbf{FALSIFIED/RETIRED:} every individual scale-one edge necessarily prevents \(K\)-erasure in a proved typed recursive system.

\section*{STALE OR OVERSTATED CLAIMS FOUND}
The unqualified MAIN1854 statement that any theta-one channel blocks exponent erasure is too broad after typed mutual recursion is admitted. Restrict it to unresolved direct channels or use the cycle-level condition. Root main's analogous cycle statement is also too broad. No load-bearing OPEN arrow closes.

\section*{Single next target}
For the first actual source-matched return construction, emit the complete typed channel graph with
\[
(\eta_e,\theta_e,\text{actual length rule},\text{workspace},\text{SAME\_REALIZATION witness}).
\]
Compute strongly connected components and cycle gains. For each same-scale edge prove well-foundedness and audit whether its query can be streamed without linear explicit workspace. The missing load-bearing theorem remains source acquisition from an arbitrary hypothetical polynomial-time exact SAT decider.

\section*{References}
S. R. Buss and R. R. Williams, \emph{Limits on Alternation-Trading Proofs for Time-Space Lower Bounds}, Computational Complexity 24(3), 2015.
L. Fortnow, R. Lipton, D. van Melkebeek, and A. Viglas, \emph{Time-Space Lower Bounds for Satisfiability}, Journal of the ACM 52(6), 2005.

\section*{FINAL STATUS}
CYCLEGAIN1 is \textbf{PROVED/EXACT conditional deterministic recurrence accounting} under its declared finite typed graph, cycle-contraction, nonexpanding same-scale, well-foundedness, and charge hypotheses. It weakens the per-edge fractional condition for time closure but exposes a separate linear-space gate for materialized same-scale queries. Source-matched graph construction, arbitrary-solver acquisition, low-space realization, the nonrelativizing/nonalgebrizing hinge, and final separation remain OPEN. P versus NP remains OPEN.

% END_OF_VERIFIED_REPORT::LSPV-3::CYCLEGAIN1
\end{document}
